\documentclass[11pt]{article}
\usepackage[a4paper,margin=1in]{geometry}
\usepackage[T1]{fontenc}
\usepackage{array}
\usepackage{hyperref}
\usepackage{longtable}
\usepackage{enumitem}

\setlist[itemize]{leftmargin=1.5em}
\setlist[enumerate]{leftmargin=1.5em}

\title{SHAKTI C-Class Performance Model Notes}
\author{perf-model}
\date{\today}

\begin{document}
\maketitle

\section{Goal}

The model is a Python cycle-level performance model for the single-issue SHAKTI
C-Class RISC-V core. It follows the CVA6 performance-model methodology: simulate
performance, not architectural behavior. The model consumes the dynamic
instruction stream that the RTL actually committed, then predicts the commit
cycle for each instruction.

The current calibrated CoreMark result is:

\begin{center}
\begin{tabular}{|l|r|}
\hline
Metric & Value \\
\hline
CoreMark iterations & 40 \\
RTL cycles from \texttt{app\_log} & 14,965,227 \\
Model cycles & 14,854,777 \\
Cycle delta & -110,450 (-0.738\%) \\
Inter-commit delta accuracy & 99.233944\% \\
Matched intervals & 12,666,565 / 12,764,347 \\
\hline
\end{tabular}
\end{center}

The held-out Dhrystone result after CoreMark calibration is 99.655836\%
inter-commit delta accuracy, with model cycles 167,113 versus 167,827 RTL cycles.

\section{RTL Trace Instrumentation}

The original C-Class commit log did not include cycle numbers. Cycle accuracy
requires a timestamp on each committed instruction, so a minimal testbench-only
change was made to \texttt{test\_soc/c64\_c32/TbSoc.bsv}. The core datapath was
not modified.

Each commit line is emitted as:

\begin{verbatim}
cycle <cycle_count> core   0: ...
\end{verbatim}

This produces the RTL commit cycle for every dynamic instruction. Benchmark-level
cycle and instruction counts were checked after instrumentation to confirm that
the tracing change was non-perturbing.

\section{Trace Parser}

\texttt{trace.py} parses the RTL \texttt{rtldump} commit log into
\texttt{TraceEntry} objects. Each entry records:

\begin{itemize}
  \item dynamic instruction index;
  \item program counter;
  \item instruction encoding;
  \item RTL commit cycle;
  \item privilege mode;
  \item integer/floating register writes;
  \item CSR writes;
  \item memory addresses when present;
  \item inferred actual next PC.
\end{itemize}

The actual next PC is inferred from the next committed instruction. That is enough
to classify branches as taken or fallthrough without computing data values.

\section{Timing Decoder}

\texttt{isa.py} decodes RV64IMAFDC only far enough for timing. The model does not
need to know that an add computes a sum, but it does need to know which registers
are read and written, which functional unit is used, whether an instruction is a
load, store, branch, jump, multiply, divide, CSR instruction, or trap, and whether
the instruction is 16-bit compressed or 32-bit.

For example, the timing metadata for a load must say:

\begin{verbatim}
reads x<base>
writes x<rd>
is_load = true
functional unit = memory
result is not immediately bypassable from execute
\end{verbatim}

This metadata drives scoreboard locking, RAW hazard checks, bypass decisions,
functional-unit latency, and front-end control timing.

Compressed decode was a concrete accuracy issue. RV64C integer doubleword
instructions such as \texttt{c.ld}, \texttt{c.sd}, \texttt{c.ldsp}, and
\texttt{c.sdsp} must not be confused with floating compressed memory operations
such as \texttt{c.fld}, \texttt{c.fsd}, \texttt{c.fldsp}, and \texttt{c.fsdsp}.
If an integer load is decoded as a floating load, the model locks the wrong
register file and misses real integer RAW dependencies.

Example:

\begin{verbatim}
c.ld   x10, 0(x8)
addi   x11, x10, 1
\end{verbatim}

Correct decode sees a dependency from \texttt{x10} to the following
\texttt{addi}. Incorrect decode as \texttt{c.fld} would mark the producer as
writing \texttt{f10}, so the model would issue the \texttt{addi} too early.

\section{Pipeline Model}

The SHAKTI C-Class is not modeled as CVA6. The useful CVA6 lesson is the
simulation methodology, not the microarchitecture. C-Class is strictly in-order
and has six major stages:

\begin{center}
\begin{tabular}{|c|l|}
\hline
Stage & Role \\
\hline
0 & PC generation and branch prediction \\
1 & instruction fetch response and compressed decompression \\
2 & decode and register read \\
3 & execute \\
4 & memory and late execute routing \\
5 & writeback, commit, CSR, and trap handling \\
\hline
\end{tabular}
\end{center}

The RTL uses inter-stage buffers (ISBs). Each stage fires when its input is not
empty, its output is not full, and its local rule guards are true. This means
backpressure is structural: a full downstream queue stalls the producer.

The model represents ISBs as bounded queues:

\begin{verbatim}
q_s0s1, q_s1s2, q_s2s3, q_s3s4, q_s4s5
\end{verbatim}

Their depths are derived from \texttt{makefile.inc}, not guessed. This matters
because changing an ISB depth changes when a local stall propagates backward
through the front end.

One model cycle executes stages in reverse order:

\begin{verbatim}
commit
stage4
execute
decode
fetch/decompress
fetch
\end{verbatim}

This approximates synchronous hardware behavior. Downstream stages consume first,
then upstream stages can observe freed queue space in the same modeled cycle. A
forward-order software loop would create artificial stalls.

\section{Modeled Hazards}

The model explicitly includes:

\begin{itemize}
  \item structural hazards from full ISBs;
  \item in-order issue and commit;
  \item scoreboard register locks;
  \item scoreboard release at commit;
  \item bypass from the heads of the two later ISBs;
  \item fixed load/store cache-hit latency;
  \item multiplier and divider latency from \texttt{MULSTAGES\_TOTAL} and
        \texttt{DIVSTAGES};
  \item branch misprediction refill penalty;
  \item CSR/trap/fence-style writeback flush penalty.
\end{itemize}

The scoreboard is modeled as a per-register lock array, matching C-Class rather
than CVA6's out-of-order scoreboard FIFO. Destination registers are locked at
execute and released at commit. With \texttt{no\_wawstalls}, the model tracks a
small WAW identifier so a late release does not unlock a newer writer.

\section{Branch Predictor}

Most residual timing errors were control-flow related, so the branch predictor
had to be modeled in detail. The predictor model includes:

\begin{itemize}
  \item BTB entries;
  \item BHT saturating counters;
  \item gshare global history;
  \item return-address stack when RTL is built with \texttt{bpu\_ras};
  \item compressed lower/upper halfword matching through the BTB \texttt{hi} bit;
  \item delayed predictor training;
  \item global-history restore after misprediction.
\end{itemize}

The largest CoreMark accuracy jumps came from:

\begin{itemize}
  \item fixing compressed integer memory decode;
  \item matching BTB \texttt{hi} during prediction;
  \item reproducing BSV hash width behavior for gshare;
  \item applying predictor training one model cycle after execute;
  \item deriving RAS enablement from \texttt{makefile.inc}.
\end{itemize}

The gshare hash fix was subtle. In BSV, shifting a \texttt{Bit\#(histbits)}
value keeps that narrow width before widening. Python integers are unbounded, so
the model had to explicitly mask the shifted history value to reproduce RTL
indexing.

\section{Accuracy Metric}

The model uses the CVA6 inter-commit spacing metric. For dynamic instruction
\(i\):

\[
\Delta t_{\mathrm{RTL}}(i) = t_{\mathrm{RTL}}(i) - t_{\mathrm{RTL}}(i-1)
\]

\[
\Delta t_{\mathrm{model}}(i) = t_{\mathrm{model}}(i) - t_{\mathrm{model}}(i-1)
\]

\[
\mathrm{accuracy} =
\frac{|\{i : \Delta t_{\mathrm{model}}(i) =
\Delta t_{\mathrm{RTL}}(i)\}|}{|\{i\}|}
\]

This metric does not let one early offset hide later local timing mistakes.

\section{What Did Not Help}

Several plausible fixes were tested and rejected:

\begin{itemize}
  \item disabling RAS reduced CoreMark accuracy to roughly 99.209439\%;
  \item adding a global extra cycle to mispredicts reduced accuracy to 98.664530\%;
  \item a coarse wrong-path front-end occupancy model also reduced accuracy to
        98.664530\%.
\end{itemize}

These negative results show why the current model should not be tuned with broad
constants. The remaining mismatches are conditional front-end/control effects.

\section{Dual-Issue Prediction Mode}

The dual-issue mode now models the actual SHAKTI dual-issue branch rather than
an ideal independent-lane machine. The important parameters extracted from the
dual branch are:

\begin{itemize}
  \item \texttt{num\_issue=2};
  \item 64-bit fetch path;
  \item \texttt{instr\_queue=6} for the stage1-to-stage2 MIMO queue, not the
        apparent \texttt{isb\_s1s2=2};
  \item one vector bundle slot for stage2-to-stage3, represented internally as
        two scalar trace entries;
  \item eight vector bundle slots for stage3-to-stage4 and stage4-to-stage5,
        represented internally as sixteen scalar trace entries;
  \item two ALUs;
  \item one shared branch/control unit;
  \item one memory issue path;
  \item one mul/div path;
  \item one FPU path;
  \item no intra-bundle forwarding;
  \item \texttt{dual\_mem} disabled and effectively unwired.
\end{itemize}

Pairing is not arbitrary. The model implements the stage2 whitelist:

\begin{itemize}
  \item ALU may pair with ALU, MULDIV, FLOAT, MEMORY, or CONTROL;
  \item CONTROL may pair with MULDIV, FLOAT, or MEMORY;
  \item MEM+MEM is disabled because \texttt{dual\_mem} is not compiled;
  \item MULDIV+MULDIV, FLOAT+FLOAT, MEM+MULDIV, MEM+FLOAT,
        MULDIV+FLOAT, BRANCH+BRANCH, and anything with SYSTEM/TRAP/WFI are
        single-issued.
\end{itemize}

The model also implements the non-obvious pair hazards. Intra-bundle RAW from
the first fetched instruction to the second fetched instruction blocks pairing.
WAW and WAR pairs are allowed because \texttt{no\_wawstalls} is defined. Once a
pair is formed, it advances in lockstep: if either slot's operands, functional
unit, or stage4 result are not ready, neither slot advances. Stage5 commits the
pair atomically, so both instructions receive the same predicted commit cycle.

These dual-issue outputs are predictions, not validated accuracy results,
because the current input trace is still the single-issue committed instruction
stream. A dual RTL trace is needed before an inter-commit accuracy number can be
reported for the dual branch.

\section{Known Limitations}

The model deliberately does not capture:

\begin{itemize}
  \item data-cache misses or cache conflicts;
  \item instruction-cache misses;
  \item data-dependent divider timing;
  \item CSR/trap behavior requiring data values;
  \item fully accurate wrong-path front-end queue occupancy after control
        redirects.
\end{itemize}

\end{document}
